% ---------------------------------------------------------------
% preamble.tex — ตั้งค่าภาษาไทยสำหรับ XeLaTeX
% ---------------------------------------------------------------

% 1) ตัดคำไทยด้วยพจนานุกรม ICU ภายใน XeTeX
%    (ภาษาไทยไม่มีช่องว่างระหว่างคำ ถ้าไม่ตั้งค่านี้บรรทัดจะไม่ตัดหรือล้นขอบ)
\XeTeXlinebreaklocale "th"
\XeTeXlinebreakskip = 0pt plus 1pt

% 2) ป้ายต่าง ๆ ของเล่มเป็นภาษาไทย
\AtBeginDocument{%
  \renewcommand{\contentsname}{สารบัญ}%
  \renewcommand{\listfigurename}{สารบัญรูป}%
  \renewcommand{\listtablename}{สารบัญตาราง}%
  \renewcommand{\bibname}{เอกสารอ้างอิง}%
  \renewcommand{\appendixname}{ภาคผนวก}%
  \renewcommand{\partname}{ภาค}%
  \renewcommand{\chaptername}{บทที่}%
}

% 3) หัวบทแบบ "บทที่ 1" ของ KOMA-Script
\renewcommand*{\chapterformat}{\chaptername~\thechapter\autodot\enskip}
\renewcommand*{\autodot}{}

% 4) ตัวหนาในหัวข้อ: ไม่ใช้ sans เพื่อให้ฟอนต์ไทยกลมกลืนทั้งเล่ม
\setkomafont{disposition}{\normalfont\bfseries}

% 5) ย่อหน้าบรรทัดแรก 1.5em และเว้นระยะเล็กน้อยระหว่างย่อหน้า
%    ปรับตัวเลขตามรสนิยมได้เลย
\setlength{\parindent}{1.5em}
\setlength{\parskip}{0.3em plus 0.1em}

% 6) ตัวไทยในโค้ดและ verbatim: ฟอนต์ mono (DejaVu Sans Mono) ไม่มีตัวไทย
%    ใช้ ucharclasses สลับไป Google Sans เฉพาะช่วงตัวอักษรไทย และเปิดใช้เฉพาะใน \ttfamily
%    (เพิ่มในเล่มนี้ 23 พ.ย. 2026 ไม่มีใน Plant_Epi_Thai_book)
\usepackage{ucharclasses}
\newfontfamily\thaimonofallback[Path=./fonts/, Extension=.ttf,
  UprightFont=*-Regular, BoldFont=*-Bold, ItalicFont=*-Italic,
  BoldItalicFont=*-Bold]{GoogleSans}
\setTransitionsFor{Thai}{\begingroup\thaimonofallback}{\endgroup}
\AtBeginDocument{\XeTeXinterchartokenstate=0 }
\AddToHook{cmd/ttfamily/after}{\XeTeXinterchartokenstate=1 }
